\documentclass[11pt]{article}
\usepackage[margin=1in]{geometry}
\usepackage{amsmath,amssymb}
\usepackage{enumitem}
\setlist{nosep,leftmargin=1.3em}
\setlength{\parskip}{4pt}
\setlength{\parindent}{0pt}

\begin{document}

\begin{center}
{\large\bfseries Project Roadmap: Fair Influence Blocking Under Unequal Receptiveness}\\[2pt]
Rachel Kan and Matthew Zhang
\end{center}

Our project continues the exploratory paper. Its findings (saved-node fairness is
nearly free because of shielding, converted-node fairness is costly, strict parity
levels down while maximin does not, and the ceiling model is a poor PoF benchmark)
all come from a single 500-node SBM with one parameter setting. The project has
four future goals: test whether those findings generalize, repair the analytic model
where it failed, compare maximin and parity against other fairness notions from the
literature, and measure how far our greedy heuristics are from optimal.

\textbf{1. Robustness and real networks (main computational work).}
\begin{itemize}
  \item Rerun the core experiments over $\sim$20 independent SBM draws and report
        confidence intervals, so we can tell which effects are larger than noise.
  \item Sweep the variables we held fixed: homophily ($p_{\mathrm{out}}/p_{\mathrm{in}}$),
        group-size ratio, acceptance spread $\sigma$, IC probability and number of
        rumor seeds, and extend the parity $\lambda$ sweep to other $\mu_L$ and $k$.
        Our main prediction: since shielding needs cross-group edges, saved-node
        fairness should stop being free as homophily increases.
  \item Real graphs with natural group labels: the Political Blogs network
        (Adamic--Glance, $\sim$1.2k nodes, two homophilous groups) as the main case,
        and a network with many labeled communities (e.g.\ SNAP email-Eu-core,
        departments) for the maximin-fragility question. Acceptance will still be
        assigned per group, so these are semi-synthetic, and we will say so.
\end{itemize}

\textbf{2. A better analytic model.} Replace the single yield $b$ with a yield
matrix $\beta_{gh}$ (truth converted in group $h$ per seed in group $g$), which
captures cross-group spillover, and derive the resulting maximin and parity
allocations and PoF. We will check whether a matrix calibrated from first-seed
yields predicts the simulated allocations better than the original model.

\textbf{3. Other fairness notions.} We will compare maximin and parity against two
more objectives from our readings, in simulation and in the ceiling model.
\emph{Welfare} ($W_\alpha$, Rahmattalabi et al.) spans total blocking to leximin with
one parameter, so we can trace the full PoF--worst-group trade-off and check whether
some $\alpha$ beats both maximin and loose parity. \emph{Diversity constraints} (DC)
guarantee each group what a size-proportional share of the budget would get it, but
do not compensate for low receptiveness; we expect DC to be much cheaper than maximin
on converted nodes while leaving a larger gap.

\textbf{4. Gap to optimal.} Rejection makes even total blocking non-submodular, so
greedy has no guarantee. On small graphs (Karate Club, $\sim$30--40-node SBMs) we
will brute-force the optimum for small $k$ and measure the gap for total greedy,
worst-first maximin, DP greedy and the new objectives.

\textbf{Plan and division of work.} We keep the exploratory split (Rachel: parity,
Matthew: maximin) and share the infrastructure.
\begin{itemize}
  \item \emph{By Nov 23 (slides):} multi-draw runs with error bars, homophily sweep,
        Political Blogs loaded. Planned preview: how the cost of fairness on saved
        nodes changes with homophily.
  \item \emph{By Nov 30 (before the one-on-one check-in):} remaining sweeps,
        real-network results, yield-matrix model, welfare and DC (simulation and
        ceiling model), and the brute-force comparison. All experiments and analysis
        finished.
  \item \emph{Dec 1--15:} draw conclusions from results and write the final report.
\end{itemize}
Our priorities, in order, are (1), (3), (2) and (4); if time runs short we will
cut the brute-force comparison first.

\end{document}
